% Options for packages loaded elsewhere
% Options for packages loaded elsewhere
\PassOptionsToPackage{unicode}{hyperref}
\PassOptionsToPackage{hyphens}{url}
\PassOptionsToPackage{dvipsnames,svgnames,x11names}{xcolor}
%
\documentclass[
  letterpaper,
  DIV=11,
  numbers=noendperiod]{scrartcl}
\usepackage{xcolor}
\usepackage{amsmath,amssymb}
\setcounter{secnumdepth}{-\maxdimen} % remove section numbering
\usepackage{iftex}
\ifPDFTeX
  \usepackage[T1]{fontenc}
  \usepackage[utf8]{inputenc}
  \usepackage{textcomp} % provide euro and other symbols
\else % if luatex or xetex
  \usepackage{unicode-math} % this also loads fontspec
  \defaultfontfeatures{Scale=MatchLowercase}
  \defaultfontfeatures[\rmfamily]{Ligatures=TeX,Scale=1}
\fi
\usepackage{lmodern}
\ifPDFTeX\else
  % xetex/luatex font selection
\fi
% Use upquote if available, for straight quotes in verbatim environments
\IfFileExists{upquote.sty}{\usepackage{upquote}}{}
\IfFileExists{microtype.sty}{% use microtype if available
  \usepackage[]{microtype}
  \UseMicrotypeSet[protrusion]{basicmath} % disable protrusion for tt fonts
}{}
\makeatletter
\@ifundefined{KOMAClassName}{% if non-KOMA class
  \IfFileExists{parskip.sty}{%
    \usepackage{parskip}
  }{% else
    \setlength{\parindent}{0pt}
    \setlength{\parskip}{6pt plus 2pt minus 1pt}}
}{% if KOMA class
  \KOMAoptions{parskip=half}}
\makeatother
% Make \paragraph and \subparagraph free-standing
\makeatletter
\ifx\paragraph\undefined\else
  \let\oldparagraph\paragraph
  \renewcommand{\paragraph}{
    \@ifstar
      \xxxParagraphStar
      \xxxParagraphNoStar
  }
  \newcommand{\xxxParagraphStar}[1]{\oldparagraph*{#1}\mbox{}}
  \newcommand{\xxxParagraphNoStar}[1]{\oldparagraph{#1}\mbox{}}
\fi
\ifx\subparagraph\undefined\else
  \let\oldsubparagraph\subparagraph
  \renewcommand{\subparagraph}{
    \@ifstar
      \xxxSubParagraphStar
      \xxxSubParagraphNoStar
  }
  \newcommand{\xxxSubParagraphStar}[1]{\oldsubparagraph*{#1}\mbox{}}
  \newcommand{\xxxSubParagraphNoStar}[1]{\oldsubparagraph{#1}\mbox{}}
\fi
\makeatother


\usepackage{longtable,booktabs,array}
\usepackage{calc} % for calculating minipage widths
% Correct order of tables after \paragraph or \subparagraph
\usepackage{etoolbox}
\makeatletter
\patchcmd\longtable{\par}{\if@noskipsec\mbox{}\fi\par}{}{}
\makeatother
% Allow footnotes in longtable head/foot
\IfFileExists{footnotehyper.sty}{\usepackage{footnotehyper}}{\usepackage{footnote}}
\makesavenoteenv{longtable}
\usepackage{graphicx}
\makeatletter
\newsavebox\pandoc@box
\newcommand*\pandocbounded[1]{% scales image to fit in text height/width
  \sbox\pandoc@box{#1}%
  \Gscale@div\@tempa{\textheight}{\dimexpr\ht\pandoc@box+\dp\pandoc@box\relax}%
  \Gscale@div\@tempb{\linewidth}{\wd\pandoc@box}%
  \ifdim\@tempb\p@<\@tempa\p@\let\@tempa\@tempb\fi% select the smaller of both
  \ifdim\@tempa\p@<\p@\scalebox{\@tempa}{\usebox\pandoc@box}%
  \else\usebox{\pandoc@box}%
  \fi%
}
% Set default figure placement to htbp
\def\fps@figure{htbp}
\makeatother


% definitions for citeproc citations
\NewDocumentCommand\citeproctext{}{}
\NewDocumentCommand\citeproc{mm}{%
  \begingroup\def\citeproctext{#2}\cite{#1}\endgroup}
\makeatletter
 % allow citations to break across lines
 \let\@cite@ofmt\@firstofone
 % avoid brackets around text for \cite:
 \def\@biblabel#1{}
 \def\@cite#1#2{{#1\if@tempswa , #2\fi}}
\makeatother
\newlength{\cslhangindent}
\setlength{\cslhangindent}{1.5em}
\newlength{\csllabelwidth}
\setlength{\csllabelwidth}{3em}
\newenvironment{CSLReferences}[2] % #1 hanging-indent, #2 entry-spacing
 {\begin{list}{}{%
  \setlength{\itemindent}{0pt}
  \setlength{\leftmargin}{0pt}
  \setlength{\parsep}{0pt}
  % turn on hanging indent if param 1 is 1
  \ifodd #1
   \setlength{\leftmargin}{\cslhangindent}
   \setlength{\itemindent}{-1\cslhangindent}
  \fi
  % set entry spacing
  \setlength{\itemsep}{#2\baselineskip}}}
 {\end{list}}
\usepackage{calc}
\newcommand{\CSLBlock}[1]{\hfill\break\parbox[t]{\linewidth}{\strut\ignorespaces#1\strut}}
\newcommand{\CSLLeftMargin}[1]{\parbox[t]{\csllabelwidth}{\strut#1\strut}}
\newcommand{\CSLRightInline}[1]{\parbox[t]{\linewidth - \csllabelwidth}{\strut#1\strut}}
\newcommand{\CSLIndent}[1]{\hspace{\cslhangindent}#1}



\setlength{\emergencystretch}{3em} % prevent overfull lines

\providecommand{\tightlist}{%
  \setlength{\itemsep}{0pt}\setlength{\parskip}{0pt}}



 


\KOMAoption{captions}{tableheading}
\makeatletter
\@ifpackageloaded{caption}{}{\usepackage{caption}}
\AtBeginDocument{%
\ifdefined\contentsname
  \renewcommand*\contentsname{Table of contents}
\else
  \newcommand\contentsname{Table of contents}
\fi
\ifdefined\listfigurename
  \renewcommand*\listfigurename{List of Figures}
\else
  \newcommand\listfigurename{List of Figures}
\fi
\ifdefined\listtablename
  \renewcommand*\listtablename{List of Tables}
\else
  \newcommand\listtablename{List of Tables}
\fi
\ifdefined\figurename
  \renewcommand*\figurename{Figure}
\else
  \newcommand\figurename{Figure}
\fi
\ifdefined\tablename
  \renewcommand*\tablename{Table}
\else
  \newcommand\tablename{Table}
\fi
}
\@ifpackageloaded{float}{}{\usepackage{float}}
\floatstyle{ruled}
\@ifundefined{c@chapter}{\newfloat{codelisting}{h}{lop}}{\newfloat{codelisting}{h}{lop}[chapter]}
\floatname{codelisting}{Listing}
\newcommand*\listoflistings{\listof{codelisting}{List of Listings}}
\makeatother
\makeatletter
\makeatother
\makeatletter
\@ifpackageloaded{caption}{}{\usepackage{caption}}
\@ifpackageloaded{subcaption}{}{\usepackage{subcaption}}
\makeatother
\usepackage{bookmark}
\IfFileExists{xurl.sty}{\usepackage{xurl}}{} % add URL line breaks if available
\urlstyle{same}
\hypersetup{
  pdftitle={Agency loops for students in an age of Agentic AI},
  pdfkeywords={human agency, agentic AI, CAIL, turing trap, AI Literacy,
belief},
  colorlinks=true,
  linkcolor={blue},
  filecolor={Maroon},
  citecolor={Blue},
  urlcolor={Blue},
  pdfcreator={LaTeX via pandoc}}


\title{Agency loops for students in an age of Agentic AI}
\author{Frederik J van Deventer}
\date{}
\begin{document}
\maketitle


\subsection{Introduction}\label{introduction}

\emph{``But Chat said this was a good solution''} is something I hear in
a number of variations from my students, where Chat means ChatGPT or
another chat-based LLM project. Students are learning to work with these
systems more and more (Kumar et al. 2024) as AI is playing a bigger part
in student's academic work, school projects and school policy.

The freedom to act according to one's own intention and volition is
often described as agency, or a level of control (Skinner, n.d.). It is
this freedom that is described as \emph{achieved agency} by Biesta and
Tedder (2007) when it is \textbf{realised} ``through the active
engagement of individuals with aspects of their contexts-for-action'',
that always exists in an temporal-relational \emph{ecology} (Biesta and
Tedder 2006). This \emph{achieved agency} is under pressure when systems
become more important in decision making and can even influence personal
decision-making with ``hypernudging'' (Yeung 2017). With systems with AI
responsibility circulates among human and non-human actors, creating
``agency loops'' (Leonardi 2025). Sartre (1946) said that our freedom
``depends entirely upon the freedom of others and that the freedom of
others depends upon our own''. If systems themselves are displaying
agency are they then infringing on our human agency?

Over the last decade the question whether Artificial Intelligence (AI)
itself exhibits agency has been debated (Legaspi, He, and Toyoizumi
2019; Swanepoel 2021). This distinction between human and technological
agency becomes more blurred with the emerging paradigm of Agentic AI
(AAI), which is different to previous forms of AI such as Generative AI
(GenAI), which we encounter in platforms like ChatGPT, Claude or
Co-pilot (Naik, Naik, and Naik 2024), because AAI can operate
(semi-)autonomously and pursue complex goals with little human
intervention (Acharya, Kuppan, and Divya 2025). Both GenAI and AAI
largely depend on recent developments in Large Language Models (LLMs).
With AAI multiple instances of these GenAI-bots run in the background
interfacing each other creating seemingly autonomous and agentic
systems, as a evolutionary next step from the GenAI chat-platforms
widely in use among students already (Silvennoinen, Aksovaara, and
Alanko-Turunen 2025; Dai 2025).

The \emph{iterational} aspect of Emirbayer and Mische (1998)'s agency
upon which Biesta and Tedder (2007) based the \emph{achieved agency} is
exposed in structure and collective beliefs, as truth expressed in
instruments such as policies or punitive action (Foucault 1980, 131), or
as expressed by teachers in their narratives or dispositions (Nespor
1987). Beliefs about ``inevitability of technology'' and the cult of
innovation (Winner 2018) pervade, while at the same time mistrust and
weariness tries to keeps AI out of universities (Guest et al. 2025).
Technology is never just neutral nor is it inherently good or bad
(Morrow 2014; Heyndels 2023) it is there to be assistive to our human
needs (Arora 2024, 32), and changes us just as much as we change it
(Latour 1994, 33). Regulation and policy fueled by ideology (both
benefiting the few or the many) has been the driver for the direction
this advancement would take us (Johnson and Acemoglu 2023, 57).

Students navigate these punitive faculties, policies and teacher
narratives as ecology in which they need to realize goals, their agency
changing from situation to situation. Is Agentic AI infringing on their
intentional volitions realized as \emph{achieved agency}?

\subsection{Knowledge Gap}\label{knowledge-gap}

Most literature focusing on AI is not (yet) mentioning AAI, but using
GenAI, which is closely related as precursor to AAI paradigm that also
utilizes LLMs. This firstly shows the need to further examine the
phenomenon of AAI in relation to students. Secondly where GenAI is used
in studies we can use as corollary for the type of results AAI.

The way AAI or GenAI is influencing agency in students seems not to have
been studied, apart from the influence it has on critical thinking
(Suriano et al. 2025), and why and how they use it for studying (Dai
2025; Silvennoinen, Aksovaara, and Alanko-Turunen 2025; Suonpää,
Heikkilä, and Dimkar 2024). The Human AI - Interaction framework
suggested by Sundar (2020) incorporates agency, but mostly highlights
interesting avenues to pursue rather than diving into topics such as:
human and AI collaboration, the seeming either/or thinking of machine vs
human agency; agency is seen as a ``mediating variable'' in the
machine/human interaction (Sundar 2020).

Agentic AI and AI goals in general, largely seem to focus on increasing
abilities for models or replacing tedious tasks performed by humans.
While research has been done about future projections of human agency in
conjunction with AI in decision making (Pew Research Center, Anderson,
and Rainie 2023) and in education (Mouta, Pinto-Llorente, and
Torrecilla-Sánchez 2025) it does not deal with how achieved agency for
humans now. Or how AAI is changing realized intentions.

Theorizations of agency by Biesta and Tedder (2007) clarify that action
flows from the values and beliefs of a person. This study will try to
deepen understanding on this interplay of beliefs and agency realised in
the context of Agentic AI among students.

\subsection{Research Questions}\label{research-questions}

\subsubsection{Main research question}\label{main-research-question}

How is Agentic AI influencing beliefs among students and impacting
agency as achieved?

\paragraph{Sub Questions}\label{sub-questions}

\begin{enumerate}
\def\labelenumi{\arabic{enumi}.}
\tightlist
\item
  How are values and beliefs impacting agency as achieved?
\item
  How are prevalent beliefs disclosed within the ecology of achieved
  agency changing policy and teacher narratives because of AAI?
\item
  In what ways are students performing day-to-day tasks, by either
  conforming or subverting rules and restrictions?
\end{enumerate}

\subsection{Methodology}\label{methodology}

\begin{longtable}[]{@{}
  >{\raggedright\arraybackslash}p{(\linewidth - 6\tabcolsep) * \real{0.0385}}
  >{\raggedright\arraybackslash}p{(\linewidth - 6\tabcolsep) * \real{0.7308}}
  >{\raggedright\arraybackslash}p{(\linewidth - 6\tabcolsep) * \real{0.1538}}
  >{\raggedright\arraybackslash}p{(\linewidth - 6\tabcolsep) * \real{0.0769}}@{}}
\caption{Research Questions and how they relate to
methods}\tabularnewline
\toprule\noalign{}
\begin{minipage}[b]{\linewidth}\raggedright
\end{minipage} & \begin{minipage}[b]{\linewidth}\raggedright
Question
\end{minipage} & \begin{minipage}[b]{\linewidth}\raggedright
Method
\end{minipage} & \begin{minipage}[b]{\linewidth}\raggedright
minimum number of participants
\end{minipage} \\
\midrule\noalign{}
\endfirsthead
\toprule\noalign{}
\begin{minipage}[b]{\linewidth}\raggedright
\end{minipage} & \begin{minipage}[b]{\linewidth}\raggedright
Question
\end{minipage} & \begin{minipage}[b]{\linewidth}\raggedright
Method
\end{minipage} & \begin{minipage}[b]{\linewidth}\raggedright
minimum number of participants
\end{minipage} \\
\midrule\noalign{}
\endhead
\bottomrule\noalign{}
\endlastfoot
RQ1 & How are values and beliefs impacting agency as achieved? &
Systematic Literature Review & not applicable \\
RQ2 & How are prevalent beliefs disclosed within the ecology of achieved
agency changing policy and teacher narratives because of AAI? & Critical
Discourse Analyis & n = 5-8 \\
RQ3 & In what ways are students performing day-to-day tasks, by either
conforming or subverting rules and restrictions? & Mixed-Method & n =
5-8 for interviews n \textasciitilde{} 200 for surveys \\
\end{longtable}

\subsubsection{RQ1: Systematic Literature
Review}\label{rq1-systematic-literature-review}

To find out what studies related to agency a systematic literature will
be conducted that searches for studies that relate to agency..

\subsubsection{RQ2: Critical Discourse
Analyis}\label{rq2-critical-discourse-analyis}

Through \emph{critical discourse analyis} (CDA) I will look at what kind
of language is constructing narratives, describes relations of power and
how it creates or reinforces beliefs (Blommaert and Bulcaen 2000).

CDA tries to find patterns in power dynamics and social effects in
discourse. In this analysis we can take the texts further and analyze
context and intertextuality to challenge social inequalities. I will
focus mainly on step two: \emph{Identify obstacles to addressing the
social wrong} of this process which consists of three stages (Fairclough
2013):

\begin{enumerate}
\def\labelenumi{\arabic{enumi}.}
\tightlist
\item
  Analyse dialectical relations
\item
  Select texts
\item
  Analyse texts interdiscursively and linguistically
\end{enumerate}

To assist in the Critical Discourse analysis I will use the Cui et al.
(2023) package for R to help quantify some aspects of it.

\subsubsection{RQ3: Mixed-method}\label{rq3-mixed-method}

\emph{Thematic analysis} will be used as Naeem et al. (2023) and Braun
and Clarke (2021) describe it to find common themes and topics in
personal qualitative interview with students. Thematic analysis is a
more subjective approach to data, where language is central. It requires
a \emph{coding} process that uses semantic and latent codes to
categorize statements and find \emph{themes} and \emph{subthemes} of
meaning in texts. The coding will be assisted by the Huang (2016)
package

After the initial analysis and recognized larger themes questions will
be formulated to be able to use a more quantitative approach in the form
of a survey.

\phantomsection\label{refs}
\begin{CSLReferences}{1}{0}
\bibitem[\citeproctext]{ref-acharyaAgentic2025}
Acharya, Deepak Bhaskar, Karthigeyan Kuppan, and B. Divya. 2025.
{``Agentic {AI}: {Autonomous} Intelligence for Complex Goals---a
Comprehensive Survey.''} \emph{IEEE Access : Practical Innovations, Open
Solutions} 13: 18912--36.
\url{https://doi.org/10.1109/ACCESS.2025.3532853}.

\bibitem[\citeproctext]{ref-arora2024pessimism}
Arora, Payal. 2024. \emph{From Pessimism to Promise: {Lessons} from the
{Global South} on Designing Inclusive Tech}. MIT Press.

\bibitem[\citeproctext]{ref-biestaHowAgencyPossible}
Biesta, Gert, and M. Tedder. 2006. {``How Is Agency Possible? {Towards}
an Ecological Understanding of Agency-as-Achievement.''} Working paper
5. The Learning Lives project.

\bibitem[\citeproctext]{ref-biestaAgencyLearningLifecourse2007}
Biesta, Gert, and Michael Tedder. 2007. {``Agency and Learning in the
Lifecourse: {Towards} an Ecological Perspective.''} \emph{Studies in the
Education of Adults} 39 (2): 132--49.
\url{https://doi.org/10.1080/02660830.2007.11661545}.

\bibitem[\citeproctext]{ref-blommaertCriticalDiscourseAnalysis2000}
Blommaert, Jan, and Chris Bulcaen. 2000. {``Critical {Discourse
Analysis}.''} \emph{Annual Review of Anthropology} 29: 447--66.
\url{https://www.jstor.org/stable/223428}.

\bibitem[\citeproctext]{ref-braun2021thematic}
Braun, Virginia, and Victoria Clarke. 2021. {``Thematic Analysis: A
Practical Guide to Understanding and Doing.''} \emph{Thousand Oaks}.

\bibitem[\citeproctext]{ref-Cui2023}
Cui, Qi, Joshua P. Le, Albert Chai, Andrew S. Lee Without Orcid, Jitarth
Sheth Without Orcid, Kevin Banh Without Orcid, Priya Pahal Without
Orcid, Katherine Ly Without Orcid, and Stanley M. Lo. 2023.
{``{discourseGT}: {An R} Package to Analyze Discourse Networks in
Educational Contexts.''} \emph{Journal of Open Source Software} 8 (88):
5143. \url{https://doi.org/10.21105/joss.05143}.

\bibitem[\citeproctext]{ref-daiWhyStudentsUse2025}
Dai, Yun. 2025. {``Why Students Use or Not Use Generative {AI}:
{Student} Conceptions, Concerns, and Implications for Engineering
Education.''} \emph{Digital Engineering} 4 (March): 100019.
\url{https://doi.org/10.1016/j.dte.2024.100019}.

\bibitem[\citeproctext]{ref-emirbayerWhatAgency1998}
Emirbayer, Mustafa, and Ann Mische. 1998. {``What {Is Agency}?''}
\emph{American Journal of Sociology} 103 (4): 962--1023.
\url{https://doi.org/10.1086/231294}.

\bibitem[\citeproctext]{ref-fairclough2013critical}
Fairclough, Norman. 2013. \emph{Critical Discourse Analysis: {The}
Critical Study of Language}. Routledge.

\bibitem[\citeproctext]{ref-Foucault1980TruthPower}
Foucault, Michel. 1980. {``Truth and Power.''} In \emph{Power/Knowledge:
{Selected} Interviews and Other Writings, 1972--1977}, edited by Colin
Gordon, 109--33. New York: Pantheon Books.

\bibitem[\citeproctext]{ref-guestUncriticalAdoptionAI2025}
Guest, Olivia, Marcela Suarez, Barbara Müller, Edwin van Meerkerk,
Arnoud Oude Groote Beverborg, Ronald de Haan, Andrea Reyes Elizondo, et
al. 2025. {``Against the {Uncritical Adoption} of '{AI}' {Technologies}
in {Academia}.''} Zenodo. \url{https://doi.org/10.5281/ZENODO.17065099}.

\bibitem[\citeproctext]{ref-heyndelsTechnologyNeutrality2023}
Heyndels, Sybren. 2023. {``Technology and {Neutrality}.''}
\emph{Philosophy \& Technology} 36 (4): 75.
\url{https://doi.org/10.1007/s13347-023-00672-1}.

\bibitem[\citeproctext]{ref-rqda}
Huang, Ronggui. 2016. {``{RQDA}: {R-based Qualitative Data Analysis}.''}

\bibitem[\citeproctext]{ref-johnsonPowerProgressOur2023}
Johnson, Simon, and Daron Acemoglu. 2023. \emph{Power and {Progress}:
{Our Thousand-Year Struggle Over Technology} and {Prosperity}}. Hachette
UK.

\bibitem[\citeproctext]{ref-kumarGuidingStudentsUsing2024}
Kumar, Harsh, Ilya Musabirov, Mohi Reza, Jiakai Shi, Xinyuan Wang,
Joseph Jay Williams, Anastasia Kuzminykh, and Michael Liut. 2024.
{``Guiding {Students} in {Using LLMs} in {Supported Learning
Environments}: {Effects} on {Interaction Dynamics}, {Learner
Performance}, {Confidence}, and {Trust}.''} \emph{Proc. ACM Hum.-Comput.
Interact.} 8 (CSCW2): 499:1--30. \url{https://doi.org/10.1145/3687038}.

\bibitem[\citeproctext]{ref-bruno1994technical}
Latour, Bruno. 1994. {``On Technical Mediation. {Philosophy}, Sociology,
Genealogy.''} \emph{Common Knowledge} 3: 29--64.

\bibitem[\citeproctext]{ref-legaspiSyntheticAgencySense2019}
Legaspi, Roberto, Zhengqi He, and Taro Toyoizumi. 2019. {``Synthetic
Agency: Sense of Agency in Artificial Intelligence.''} \emph{Current
Opinion in Behavioral Sciences}, Artificial {Intelligence}, 29
(October): 84--90. \url{https://doi.org/10.1016/j.cobeha.2019.04.004}.

\bibitem[\citeproctext]{ref-leonardiHomoAgenticusAge2025}
Leonardi, Paul M. 2025. {``{\emph{Homo Agenticus}} in the Age of Agentic
{AI}: {Agency} Loops, Power Displacement, and the Circulation of
Responsibility.''} \emph{Information and Organization} 35 (3): 100582.
\url{https://doi.org/10.1016/j.infoandorg.2025.100582}.

\bibitem[\citeproctext]{ref-morrowWhenTechnologiesMakes2014}
Morrow, David R. 2014. {``When {Technologies Makes Good People Do Bad
Things}: {Another Argument Against} the {Value-Neutrality} of
{Technologies}.''} \emph{Science and Engineering Ethics} 20 (2):
329--43. \url{https://doi.org/10.1007/s11948-013-9464-1}.

\bibitem[\citeproctext]{ref-moutaWhereAgencyMoving2025}
Mouta, Ana, Ana María Pinto-Llorente, and Eva María Torrecilla-Sánchez.
2025. {``{`{Where} Is {Agency Moving} to?'}: {Exploring} the {Interplay}
Between {AI Technologies} in {Education} and {Human Agency}.''}
\emph{Digital Society} 4 (2): 49.
\url{https://doi.org/10.1007/s44206-025-00203-9}.

\bibitem[\citeproctext]{ref-naeemStepbyStepProcessThematic2023}
Naeem, Muhammad, Wilson Ozuem, Kerry Howell, and Silvia Ranfagni. 2023.
{``A {Step-by-Step Process} of {Thematic Analysis} to {Develop} a
{Conceptual Model} in {Qualitative Research}.''} \emph{International
Journal of Qualitative Methods} 22 (October): 16094069231205789.
\url{https://doi.org/10.1177/16094069231205789}.

\bibitem[\citeproctext]{ref-naikChatGPTAllYou2024}
Naik, Ishita, Dishita Naik, and Nitin Naik. 2024. {``{ChatGPT Is All You
Need}: {Untangling Its Underlying AI Models}, {Architecture}, {Training
Procedure}, {Capabilities}, {Limitations And Applications}.''}
Preprints.
\url{https://doi.org/10.36227/techrxiv.173273427.76836200/v1}.

\bibitem[\citeproctext]{ref-nesporRoleBeliefsPractice1987}
Nespor, Jan. 1987. {``The Role of Beliefs in the Practice of
Teaching.''} \emph{Journal of Curriculum Studies} 19 (4): 317--28.
\url{https://doi.org/10.1080/0022027870190403}.

\bibitem[\citeproctext]{ref-pewFutureHuman2023}
Pew Research Center, Janna Anderson, and Lee Rainie. 2023. {``The
{Future} of {Human Agency},''} February.

\bibitem[\citeproctext]{ref-ExistentialismHumanismJeanPaul}
Sartre, Jean Paul. 1946. {``Existentialism Is a {Humanism}.''}
https://www.marxists.org/reference/archive/sartre/works/exist/sartre.htm.

\bibitem[\citeproctext]{ref-silvennoinenStudentsUsageGenAI2025}
Silvennoinen, Minna, Satu Aksovaara, and Merja Alanko-Turunen. 2025.
{``Students' {Usage} of {GenAI} in {Universities} of {Applied Sciences}:
{Experiences} and {Development Needs} for {Guidance} and {Support}.''}
In \emph{38th {Bled eConference}: {Empowering Transformation}: {Shaping
Digital Futures} for {All}: {Conference Proceedings}}, 467--82.
University of Maribor Press.
\url{https://doi.org/10.18690/um.fov.4.2025.29}.

\bibitem[\citeproctext]{ref-skinnerGuideConstructsControl}
Skinner, Ellen A. n.d. {``A {Guide} to {Constructs} of {Control}.''}

\bibitem[\citeproctext]{ref-sundarRiseMachineAgency2020}
Sundar, S Shyam. 2020. {``Rise of {Machine Agency}: {A Framework} for
{Studying} the {Psychology} of {Human}--{AI Interaction} ({HAII}).''}
\emph{Journal of Computer-Mediated Communication} 25 (1): 74--88.
\url{https://doi.org/10.1093/jcmc/zmz026}.

\bibitem[\citeproctext]{ref-suonpaaStudentsPerceptionsGenerative2024}
Suonpää, Maija, Jutta Heikkilä, and Ana Dimkar. 2024. {``Students'
{Perceptions Of Generative Ai Usage And Risks In A Finnish Higher
Education Institution}.''} In \emph{18th {International Technology},
{Education} and {Development Conference}}, 3071--77. Valencia, Spain.
\url{https://doi.org/10.21125/inted.2024.0825}.

\bibitem[\citeproctext]{ref-surianoStudentInteractionChatGPT2025}
Suriano, Rossella, Alessio Plebe, Alessandro Acciai, and Rosa Angela
Fabio. 2025. {``Student Interaction with {ChatGPT} Can Promote Complex
Critical Thinking Skills.''} \emph{Learning and Instruction} 95
(February): 102011.
\url{https://doi.org/10.1016/j.learninstruc.2024.102011}.

\bibitem[\citeproctext]{ref-swanepoelDoesArtificialIntelligence2021}
Swanepoel, Danielle. 2021. {``Does {Artificial Intelligence Have
Agency}?''} In \emph{The {Mind-Technology Problem}}, edited by Robert W.
Clowes, Klaus Gärtner, and Inês Hipólito, 18:83--104. Cham: Springer
International Publishing.
\url{https://doi.org/10.1007/978-3-030-72644-7_4}.

\bibitem[\citeproctext]{ref-winnerCultInnovationIts2018}
Winner, Langdon. 2018. {``The {Cult} of {Innovation}: {Its Myths} and
{Rituals}.''} In \emph{Engineering a {Better Future}: {Interplay}
Between {Engineering}, {Social Sciences}, and {Innovation}}, edited by
Eswaran Subrahmanian, Toluwalogo Odumosu, and Jeffrey Y. Tsao, 61--73.
Cham: Springer International Publishing.
\url{https://doi.org/10.1007/978-3-319-91134-2_8}.

\bibitem[\citeproctext]{ref-yeungHypernudgeBigData2017}
Yeung, Karen. 2017. {``{`{Hypernudge}'}: {Big Data} as a Mode of
Regulation by Design.''} \emph{Information, Communication \& Society} 20
(1): 118--36. \url{https://doi.org/10.1080/1369118X.2016.1186713}.

\end{CSLReferences}




\end{document}
